\begin{table}[htbp]
\centering
\caption{Sintesis lintas-lapisan: kapan RAG \emph{terkondisi-prediksi} memberi nilai tambah dan kapan tidak. Tabel ini merangkum temuan seluruh subbab evaluasi, termasuk yang tidak menguntungkan, agar batas manfaat metode ini terbaca jelas.}
\label{tab:sintesis}
\small
\begin{tabular}{@{}llll@{}}
\toprule
\textbf{Lapisan / skenario} & \textbf{Nilai tambah} & \textbf{Bukti} & \textbf{Sebab} \\
\midrule
Penelusuran, kasus divergen & \textbf{Besar} & MRR 0,065 $\rightarrow$ 0,168 & Kondisi terkini menyesatkan; \\
(12,9\% jendela) & \emph{dan robust} & unggul 6/6 \emph{fold} ($p=0{,}031$) & hanya prediksi yang menolong \\
\addlinespace
Penelusuran, distribusi natural & \textbf{Tidak ada} & MRR 0,731 $\rightarrow$ 0,742 & 87,1\% jendela tidak divergen: \\
& & ($p=0{,}225$; t.s.) & RAG standar sudah tepat sasaran \\
\addlinespace
Penelusuran, pengklasifikasi & \textbf{Besar} & MRR 0,225 lawan 0,168 & Pengklasifikasi tidak menyusut \\
(divergen, vs regresi) & \emph{dan robust} & unggul 6/6 \emph{fold} ($p=0{,}031$) & ke tengah sebaran \\
\addlinespace
Penelusuran, pengklasifikasi & \textbf{Merugikan} & MRR 0,698 lawan 0,742 & Kondisi terkini sudah benar; \\
(natural, vs regresi) & & kalah 6/6 \emph{fold} ($p=0{,}031$) & berpindah kelas justru menjauh \\
\addlinespace
Penelusuran, cara leksikal & \textbf{Besar} & \emph{context recall} 0,800 & Model \emph{embedding} berbahasa \\
(vs padat) & & lawan 0,200 & Inggris, korpus Indonesia \\
\addlinespace
Penelusuran, penggabungan & \textbf{Merugikan} & \emph{context recall} 0,550 & Fusi peringkat tetap memberi \\
hibrida (vs leksikal saja) & & lawan 0,800 & bobot pada daftar yang keliru \\
\midrule
Deteksi hipoglikemia & \textbf{Besar} & Sensitivitas 17,3\% $\rightarrow$ 67,2\% & Kondisi diprediksi langsung, \\
(pengklasifikasi kondisi) & & pada ambang standar & bukan diambangkan dari regresi \\
\midrule
Generasi, kasus hipoglikemia & Moderat & \emph{sim\_ref} 0,450 $\rightarrow$ 0,629 & Antisipasi mengubah tindakan \\
& & (kasus D1) & secara kualitatif \\
\addlinespace
Generasi, kasus hiperglikemia & \textbf{Tidak ada} & \emph{sim\_ref} 0,542 $\rightarrow$ 0,525 & Anjuran pedoman serupa untuk \\
& & (kasus D6) & kondisi kini maupun mendatang \\
\midrule
Prediksi SMBG, $+30$ menit & \textbf{Merugikan} & RMSE 28,56 lawan \emph{persistence} & Jarak antar-pembacaan (124,8 mnt) \\
& & 27,33~mg/dL & lebih panjang dari horizon \\
\addlinespace
Prediksi SMBG, $+60$ menit & Moderat & RMSE 46,53 lawan \emph{persistence} & Asumsi ``tidak berubah'' runtuh; \\
& & 47,87~mg/dL & efek IOB/COB mendominasi \\
\bottomrule
\end{tabular}

\vspace{0.4em}
\footnotesize t.s. = tidak signifikan. Angka penelusuran pada baris pengondisian berasal dari validasi silang enam \emph{fold} lintas-pasien (Tabel~\ref{tab:crossfold-retrieval}), sedangkan baris perbandingan cara penelusuran berasal dari metrik RAGAS berbasis rujukan (Tabel~\ref{tab:lengan-penelusuran}). Dua baris pengklasifikasi sengaja dipisah karena arahnya \textbf{berbalik} antar-himpunan; menggabungkannya menjadi satu putusan akan menyembunyikan temuan yang justru menerangkan mekanismenya. Batas atas metode ini terukur dan konsisten: dengan kondisi masa depan yang diketahui benar (\emph{oracle}), penelusuran mencapai MRR 0,833 pada kasus divergen. Seluruh selisih antara 0,168 dan 0,833 adalah ruang perbaikan yang dapat diraih semata-mata dengan meningkatkan ketepatan prediksi kondisi, tanpa mengubah apa pun pada komponen penelusuran.
\end{table}
